\documentclass{article}

\usepackage{tabularx}
\usepackage{booktabs}

\title{Problem Statement and Goals\\\progname}

\author{\authname}

\date{}

\input{../Comments.text}
\input{../Common.text}

\begin{document}

\maketitle

\begin{table}[hp]
\caption{Revision History} \label{TblRevisionHistory}
\begin{tabularx}{\textwidth}{llX}
\toprule
\textbf{Date} & \textbf{Developer(s)} & \textbf{Change}\\
\midrule
Date1 & Name(s) & Description of changes\\
Date2 & Name(s) & Description of changes\\
... & ... & ...\\
\bottomrule
\end{tabularx}
\end{table}

\section{Problem Statement}

\wss{You should check your problem statement with the
\href{https://smiths.github.io/capTemplate/Checklists/ProbState-Checklist.pdf}
{problem statement checklist}}. 

\wss{You can change the section headings, as long as you include the required
information.}

\subsection{Problem}

\wss{What problem are you trying to solve?  Why is it important? Do not talk
about how you will solve the problem, only what the problem is.  AI is rarely
part of the problem.  The problem is ``what'' you want to do, not ``how'' you
want to do it.  The problem statement should be written in a way that is
understandable to a non-technical audience.}

\subsection{Inputs and Outputs}

\subsection{Inputs}
\begin{enumerate}
    \item Taps on the screen: mobile app users will need to click the screen to complete actions in the application.
    \item Activity Information: connection to a pedometer, heart rate monitor, etc. 
        Many of these technologies can be provided by the device running the application.
    \item User Personal information: height, weight, health conditions, current fitness level, current understanding of finance.
\end{enumerate}
\subsection{Outputs}
\begin{enumerate}
    \item Virtual currency: users will be able to spend currency as they choose.
    \item Real world rewards: virtual currency will be able to trade their currency for real world incentives (i.e. free trials for paid services).
\end{enumerate}

\subsection{Stakeholders}

\subsection{Environment}

\wss{Hardware and Software Environment for the users.  The developer environment
is summarized as part of the developer plan.}

\section{Goals}

\section{Stretch Goals}

\wss{Goals you hope to get to, but not necessary.  They are also helpful if the
scope of the original project needs to be expanded.}

\section{Custom Documents (CDs)}

\wss{For CAS 741: State whether the project is a research project. This
designation, with the approval (or request) of the instructor, can be modified
over the course of the term.}

\wss{For SE Capstone: List your custom documents (CDs).  Potential CDs include
usability testing, code walkthroughs, user documentation, formal proof,
GenderMag personas, Design Thinking, etc.  (The full list is on the course
outline and in Lecture 02.) Normally the number of CDs will be two (one per
term).  Approval of the CDs will be part of the discussion with the instructor
for approving the project. The CDs, with the approval (or request) of the
instructor, can be modified over the course of the term.}

\newpage{}

\section*{Appendix --- Reflection}

\wss{Not required for CAS 741}

\input{../Reflection.text}

\begin{enumerate}
    \item What went well while writing this deliverable? 
    \item What pain points did you experience during this deliverable, and how
    did you resolve them?
    \item How did you and your team adjust the scope of your goals to ensure
    they are suitable for a Capstone project (not overly ambitious but also of
    appropriate complexity for a senior design project)?
\end{enumerate}  

\end{document}